This undergraduate thesis presents the design and implementation of the Learning to Learn and Active Listening modules for a gamified platform aimed at strengthening soft skills among Systems Engineering students at Universidad del Valle, Tuluá. The project addresses the gap between the technical emphasis of traditional engineering curricula and the increasing demand for transversal competencies essential for professional performance in contemporary work environments.
\\
\\
The development process involved a literature-based identification of theories, strategies, and pedagogical approaches related to autonomous learning, metacognition, self-regulated learning, and active listening. Based on these foundations, the platform integrates educational resources, interactive learning experiences, and gamified activities to promote motivation, participation, and meaningful practice.
\\
\\
The Learning to Learn module includes tools for goal setting, creation and management of objectives and tasks, a Kanban-based activity planner, reflective self-assessment instruments, and an AI-driven system for generating personalized learning roadmaps using natural language processing. The Active Listening module offers a game-based training experience composed of dynamic rounds with audio challenges generated through AI-based text-to-speech systems. These challenges include multiple-choice comprehension, cloze tests, paraphrasing, summarization, and clarification inquiries, each providing immediate feedback.
\\
\\
A pilot evaluation was conducted with students, using perception surveys and a heuristic-based usability assessment. Results show high levels of satisfaction, clarity of the interface, and perceived usefulness of the activities for practicing autonomous learning and active listening skills. The findings highlight the potential of gamification and AI-supported tools to foster soft-skill development in engineering education.
\\
\\
\textbf{Keywords:} Soft skills, gamification, modules, learning to learn, active listening